% from ega3-6-fr.tex
\oldpage[III]{158}
dans $Y=\operatorname{Spec}(B_1\otimes_A B_2)$ (II, 4.3.2.1); si
$X'^{(i)}=f_i^{-1}(W^{(i)})$, il s'agit d'établir un isomorphisme canonique
de foncteurs spectraux
\[
  {}^{(t)}\mathscr E
  (X'^{(1)},X'^{(2)};
   \mathscr P'_\bullet{}^{(1)},\mathscr P'_\bullet{}^{(2)})
  \xrightarrow{\sim}
  {}^{(t)}\mathscr E
  (X^{(1)},X^{(2)};
   \mathscr P_\bullet^{(1)},\mathscr P_\bullet^{(2)})
  \otimes_B B_g
  \tag{6.7.1.1}\label{III.6.7.1.1.fr}
\]
où l'on a posé $B=B_1\otimes_A B_2$ et $g=g_1\otimes g_2$. Pour cela,
partons de recouvrements ouverts affines finis $\mathfrak U^{(i)}$ de
$X^{(i)}$ ($i=1,2$), et soit $\mathfrak U'^{(i)}$ la trace de
$\mathfrak U^{(i)}$ sur $X'^{(i)}$, qui est encore formé d'ouverts affines
(I, 5.5.10); de façon précise, on a
\[
  C^\bullet(\mathfrak U'^{(i)},\mathscr P'_\bullet{}^{(i)})
  =C^\bullet(\mathfrak U^{(i)},\mathscr P_\bullet^{(i)})
   \otimes_{B_i}(B_i)_{g_i}.
\]
Si on pose
$L'_{\bullet\bullet}{}^{(i)}
=C^\bullet(\mathfrak U'^{(i)},\mathscr P'_\bullet{}^{(i)})$, on a donc
\[
  L'_{\bullet\bullet}{}^{(1)}\otimes_A L'_{\bullet\bullet}{}^{(2)}
  =\bigl(L_{\bullet\bullet}^{(1)}\otimes_{B_1}(B_1)_{g_1}\bigr)
   \otimes_A
   \bigl(L_{\bullet\bullet}^{(2)}\otimes_{B_2}(B_2)_{g_2}\bigr);
\]
comme on a $B_g=(B_1)_{g_1}\otimes_A(B_2)_{g_2}$ à un isomorphisme
canonique près, on a, à un isomorphisme canonique près,
\[
  L'_{\bullet\bullet}{}^{(1)}\otimes_A L'_{\bullet\bullet}{}^{(2)}
  =(L_{\bullet\bullet}^{(1)}\otimes_A L_{\bullet\bullet}^{(2)})
    \otimes_B B_g.
\]
Si $M_{\bullet\bullet\bullet}^{(i)}$ est une résolution projective de
Cartan-Eilenberg de $L_{\bullet\bullet}^{(i)}$, qu'on peut supposer formée de
$B_i$-modules, il résulte du fait que $(B_i)_{g_i}$ est \emph{plat} sur $B_i$
que
$M'_{\bullet\bullet\bullet}{}^{(i)}
=M_{\bullet\bullet\bullet}^{(i)}\otimes_{B_i}(B_i)_{g_i}$ est une résolution
projective de Cartan-Eilenberg du bicomplexe
$L'_{\bullet\bullet}{}^{(i)}$; en outre, on a
\[
  M'_{\bullet\bullet\bullet}{}^{(1)}\otimes_A M'_{\bullet\bullet\bullet}{}^{(2)}
  =(M_{\bullet\bullet\bullet}^{(1)}\otimes_A M_{\bullet\bullet\bullet}^{(2)})
    \otimes_B B_g.
\]
L'isomorphisme cherché (6.7.1.1) résulte alors aussitôt des définitions de
l'hyperhomologie d'un bicomplexe, et de l'exactitude du foncteur
$G\otimes_B B_g$ en le $B$-module $G$.
\end{env}

\begin{env}[6.7.2]
Supposons maintenant $S$ quelconque, et soient
$u_i:Y^{(i)}\to S$ les morphismes structuraux ($i=1,2$). Soit $(S_\alpha)$ un
recouvrement ouvert affine de $S$, posons
$Y_\alpha^{(i)}=u_i^{-1}(S_\alpha)$,
$X_\alpha^{(i)}=f_i^{-1}(Y_\alpha^{(i)})$, et soit
$f_{i\alpha}:X_\alpha^{(i)}\to Y_\alpha^{(i)}$ la restriction de $f_i$, qui
est un morphisme \emph{séparé et quasi-compact}. Les
$Y_\alpha=Y_\alpha^{(1)}\times_{S_\alpha}Y_\alpha^{(2)}$ forment un
recouvrement ouvert de $Y$, et sur chaque $Y_\alpha$ sont définis par (6.7.1)
des foncteurs spectraux
\[
  {}^{(t)}\mathscr E^{S_\alpha}
  \bigl(f_{1\alpha},f_{2\alpha};
    \mathscr P_\bullet^{(1)}|X_\alpha^{(1)},
    \mathscr P_\bullet^{(2)}|X_\alpha^{(2)}\bigr);
\]
il s'agit encore de montrer que ces foncteurs vérifient les conditions de
recollement.

On se ramène aussitôt à la situation suivante :
$S=\operatorname{Spec}(A)$ est affine, $S'=D(h)$, où $h\in A$, et
$u_i(Y^{(i)})\subset S'$; on peut en outre supposer
$Y^{(i)}=\operatorname{Spec}(B_i)$ affine; il s'agit de définir des
isomorphismes canoniques
\[
  {}^{(t)}\mathscr E^S
  (X^{(1)},X^{(2)};
   \mathscr P_\bullet^{(1)},\mathscr P_\bullet^{(2)})
  \xrightarrow{\sim}
  {}^{(t)}\mathscr E^{S'}
  (X^{(1)},X^{(2)};
   \mathscr P_\bullet^{(1)},\mathscr P_\bullet^{(2)}).
  \tag{6.7.2.1}\label{III.6.7.2.1.fr}
\]
Or, avec les notations de (6.6.2), les $L_{\bullet\bullet}^{(i)}$ sont formés
de $A_h$-modules, et l'on a donc
$L_{\bullet\bullet}^{(1)}\otimes_{A_h}L_{\bullet\bullet}^{(2)}
=L_{\bullet\bullet}^{(1)}\otimes_A L_{\bullet\bullet}^{(2)}$ à un
isomorphisme canonique près; comme on peut prendre une résolution projective
de Cartan-Eilenberg $M_{\bullet\bullet\bullet}^{(i)}$ de
$L_{\bullet\bullet}^{(i)}$ formée de $A_h$-modules, cela donne aussitôt
l'isomorphisme canonique cherché.
\end{env}

Nous avons en résumé démontré le

\begin{theorem}[6.7.3]
Soient $S$ un préschéma, $f_i:X_i\to Y_i$ un $S$-morphisme séparé et
quasi-compact de $S$-préschémas, $\mathscr P_\bullet^{(i)}$ un complexe de
$\mathscr O_{X^{(i)}}$-Modules quasi-cohérents limité inférieurement
($i=1,2$); on pose $Y=Y^{(1)}\times_S Y^{(2)}$. Il existe un
bi-$\partial$-foncteur
  $\mathfrak{Tor}_\bullet^S
(f_1,f_2;\mathscr P_\bullet^{(1)},\mathscr P_\bullet^{(2)})$
à valeurs dans la catégorie des $\mathscr O_Y$-Modules quasi-cohérents, tel
que si $V^{(i)}$ est un ouvert affine de $Y^{(i)}$ ($i=1,2$) et
$V=V^{(1)}\times_S V^{(2)}$, on ait
\[
  \mathfrak{Tor}_\bullet^S
  (f_1,f_2;\mathscr P_\bullet^{(1)},\mathscr P_\bullet^{(2)})|V
  =\Bigl(\mathbf{Tor}_\bullet^S
    \bigl(f_1^{-1}(V^{(1)}),f_2^{-1}(V^{(2)});
      \mathscr P_\bullet^{(1)}|f_1^{-1}(V^{(1)}),
      \mathscr P_\bullet^{(2)}|f_2^{-1}(V^{(2)})\bigr)\Bigr)^{\sim}.
\]
